% Abstract + MCQ "Strategy" subsection.
% Human (outline) prose is left unmarked; prose added beyond the outline is wrapped in \aiadd{}.
% Requires the preamble to define \aiadd, e.g. \newcommand{\aiadd}[1]{\textcolor{blue}{#1}}.
% Bibliography: refs/abstract_strategy.bib

\begin{abstract}
According to estimates by AI labs, benchmarking and evaluation take up a large
share of a modern training run: \aiadd{the Olmo 3 report puts evaluation at
10\% to 20\% of its compute budget}~\cite{olmo3}\aiadd{, and open efforts such
as Llama 3 run extensive benchmark suites throughout development}~\cite{llama3herd}.
Fast, granular, and trustworthy benchmarking is crucial for EDU-LLM given the
sheer number of novel hypotheses being tested. We propose a skill-based
Computerized Adaptive Test (CAT) that reduces the number of items a diagnostic
needs to predict accuracy on a specific skill, while holding a level of
precision similar to full-length benchmarks. Our findings are split between
multiple-choice benchmarks and free-response benchmarks.

On MCQ benchmarks, we reproduce ATLAS-style adaptive testing on
\aiadd{ARC}~\cite{li2025atlas}\aiadd{, where the published ATLAS item bank
recovers held-out accuracy at Pearson $r = 0.83$ from about 21 of 1,172 items
and reaches $r = 0.91$ in under 10 items once scoring protocols match}. Using
\aiadd{about 52} calibration models in the 0 to 7B range \aiadd{on benchmarks
with no public response data}, CAT recovers full benchmark accuracy with
\aiadd{1\% to 3\%} of items \aiadd{(Pedagogy $r = 0.72$, PIQA $r = 0.94$)}. On
FRQ benchmarks, we grade responses with an LLM-as-judge against per-scenario
rubrics to build a criterion-level response matrix, derive a Q-matrix, and fit
2PL/MIRT so that CAT selects scenarios adaptively.
\end{abstract}

\section{Multiple-Choice Benchmarks}

\subsection{Strategy}

\aiadd{A computerized adaptive test (CAT) chooses each question from how the
model has answered so far. It selects the item that reveals the most about the
model's ability, then stops once the estimate is precise
enough~\cite{wainer2000, vanderlinden2000cat, weiss1982}. This selection rests
on item response theory (IRT), which models the chance of a correct answer from
an item's difficulty and discrimination together with the model's latent
ability~\cite{hambleton1991}. Because the test skips items that add little, it
reaches a target precision with far fewer questions than a fixed
benchmark~\cite{maiapolo2024tiny}.}

Our work builds upon ATLAS~\cite{li2025atlas}, which showed the potential for
adaptive testing to dramatically reduce the number of items required for
predicting benchmark accuracy. Items calibrated using item response theory have
also been shown to provide information beyond
accuracy~\cite{lalor2016, rodriguez2021, vania2021}. Models that answer harder
questions correctly but miss some easy ones are appropriately ranked higher
under a CAT, even though a simple accuracy comparison cannot capture this
difference~\cite{li2025atlas}.

This paper shows the feasibility of recreating the ATLAS methodology on new
benchmarks, and it formalizes the requirements for using a CAT as a true
replacement for an MCQ benchmark. It also sets out where multidimensional IRT
(MIRT) helps and where it falls short for MCQ, a discussion largely missing
from ATLAS~\cite{li2025atlas}.
